\documentclass[12pt]{article}
\usepackage{lineno}
%\usepackage{ametsoc}
\usepackage{natbib}
\usepackage{a4wide}
\usepackage{longtable}
\usepackage{amsfonts}
\usepackage{amssymb}

\usepackage{color}
\usepackage[pdftex]{graphicx}
\usepackage[pdftex,colorlinks, citecolor=black,linkcolor=black,urlcolor=black]{hyperref}
\definecolor{darkgreen}{rgb}{0.1,0.4,0.1}
\definecolor{darkblue}{rgb}{0.1,0.1,0.3}



\newif\ifdetail
\global\detailfalse
%\global\detailtrue

\def\bem#1{\colorbox{green}{\em #1}}

\newcommand{\beq}  { \begin{eqnarray} }
\newcommand{\eeq}  { \end{eqnarray}}
\newcommand{\beeq}  { \begin{eqnarray*} }
\newcommand{\eeeq}  { \end{eqnarray*}}
\newcommand{\eq }  [1]{Eq.~(\ref{#1})}  % Eq. (#1)
\newcommand{\eqs }  [1]{Eqs.~(\ref{#1})}
\newcommand{\sect}  [1]{Sec.~(\ref{#1})}  % Eq. (#1)
\newcommand{\fig}  [1]{Fig.~\ref{#1}}   % Fig. #1
\renewcommand{\v}[1]{{\mbox{\boldmath$ #1 $}}}  % Vektor durch Fettdruck
\newcommand{\m}    [1]{{\bf  #1 }}              % Matrix in sans serif
\newcommand{\gz}[1] {{{d #1} \over {dz}}}
\newcommand{\pt}[1]{{\partial_t #1}}
\newcommand{\ptt}[1]{{\partial_{tt} #1}}
\newcommand{\pttt}[1]{{\partial_{ttt} #1}}
\newcommand{\pttx}[1]{{{\partial^3 #1} \over {\partial t^2 \partial x}}}
\newcommand{\ptty}[1]{{{\partial^3 #1} \over {\partial t^2 \partial y}}}
\newcommand{\ptx}[1]{{{\partial^2 #1} \over {\partial t \partial x}}}
\newcommand{\pty}[1]{{{\partial^2 #1} \over {\partial t \partial y}}}
\newcommand{\pz}[1]{{\partial_z #1}}
\newcommand{\pzz}[1]{{\partial_{zz} #1}}
\newcommand{\px}[1]{{\partial_x #1}}
\newcommand{\py}[1]{{\partial_y #1}}
\newcommand{\pxi}[1]{{\partial_i #1}}
\newcommand{\pxj}[1]{{\partial_j #1}}
\newcommand{\pxk}[1]{{\partial_k #1}}
\newcommand{\pb}[1]{{{\partial #1} \over {\partial b}}}
%\newcommand{\pxy}[1]{{{\partial^2 #1} \over {\partial x \partial y}}}
\newcommand{\pxy}[1]{\partial_{xy} #1}
\newcommand{\pyy}[1]{\partial_{yy} #1}
\newcommand{\pxx}[1]{\partial_{xx} #1}
\newcommand{\kx}{{\bf k} \times}
\newcommand{\dt}[1]{{{d #1} \over {d t}}}
\newcommand{\Dt}[1]{{{D #1} \over {D t}}}
\newcommand{\Div}[1]{ \nabla \cdot #1 }
\renewcommand{\div}[1]{ \nabla_{yz} \cdot #1 }
\newcommand{\Bar}[1]{ \overline{ #1} }
\newcommand{\rvec} [1]{
    \raisebox{-1.5ex}{$\stackrel{\textstyle #1}{\neg}$} }% rotierter Vektor
\newcommand{\nablq}   {\rvec\nabla}                    % rotated nabla operator
\newcommand{\Vec}[1]{\left(\begin{array}{c} #1 \end{array}\right) }
%\renewcommand{\mho}{\mathcal{W}}
\newcommand{\pn}[1]{{{\partial #1} \over {\partial n}}}
\newcommand{\pnn}[1]{{{\partial #1} \over {\partial m}}}
\newcommand{\ps}[1]{{{\partial #1} \over {\partial s}}}
\newcommand{\p}{{\partial}}
\newcommand{\vn}{{\v \nabla}}
\newcommand{\lk}{\left(}
\newcommand{\rk}{\right)}

\newcommand{\tl}{\tilde}

\newcommand{\I}{\mathcal{I}}
\newcommand{\E}{\mathcal{E}}
\newcommand{\A}{\mathcal{A}}

\newcommand{\ep}{\epsilon}
\newcommand{\erfc}{\mathrm{erfc}}
 
 
\begin{document}

 \section*{The spectral wave model wim\_realistic}

wim\_realistic is wim\_simple on the grid of pe\_realistic.
This text describes only the differences to wim\_simple.
Everything that is not mentioned here is as described in the documentation of wim\_simple.

The differences are the following.
The horizontal grid spacings may vary.
The horizontal coordinates may be longitude and latitude on the sphere.
The propagation then contains the metric factors and the turning of the rays on the sphere.
The coupling to the ocean allows a vertical grid with varying cell thickness.
The wind input of Janssen and the dissipation of WAM4 are available as an option.

Most methods of wim\_realistic are those of wim\_simple.
They are inherited or assigned from wim\_simple.
wim\_simple itself is not changed by wim\_realistic.
The source terms of wim\_simple, the nonlinear transfer, the time stepping, the reflection, the open coasts and the JAX backend are therefore the same in both models.

\subsection*{Grid}

The model domain consists of $n_x \times n_y$ cells.
The horizontal coordinates $x$ and $y$ are either Cartesian coordinates in metres or longitude $\lambda$ and latitude $\phi$ in degrees.
The attribute \verb|coord_degree| selects longitude and latitude.
The spectral grid in frequency $\nu$ and direction $\theta$ is that of wim\_simple.

The widths of the cells are given by the attributes \verb|dxt| and \verb|dyt|.
Each of them is either a number or an array of length $n_x$ or $n_y$ without halo points.
The numbers \verb|dx| and \verb|dy| of wim\_simple are accepted instead.
With \verb|coord_degree| the widths are given in degrees.

The horizontal grid is built as in pe\_realistic.
A cell has its centre at a $T$ point and its faces at $U$ points.
In $x$ the cell $i$ has the centre $x^t_i$, the width $\Delta x^t_i$ and the eastern face $x^u_i$.
The distance between the centres $x^t_i$ and $x^t_{i+1}$ is $\Delta x^u_i$.
The same holds in $y$.
The faces and centres follow from the widths by the recursion of pyOM2,
\beq \label{eq:ucentred_w}
 x^u_i = x^u_{i-1} + \Delta x^t_i ~,~~
 x^t_{i+1} = 2 x^u_i - x^t_i ~,~~
 \Delta x^u_i = x^t_{i+1} - x^t_i
\eeq
The first interior $U$ point is placed at \verb|x_origin| and \verb|y_origin| (default 0).

The recursion \eq{eq:ucentred_w} can give $\Delta x^u \le 0$ when the widths jump.
This holds also across the seam of a periodic direction.
The widths must therefore vary smoothly.

With \verb|coord_degree| the widths are converted to metres after the coordinates are built.
The factor is $a \pi /180$ with the radius of the earth $a = 6370\,$km (attribute \verb|radius|).
The metric factors are $\cos \phi$ at the $T$ rows, $\cos \phi$ at the $U$ rows and $\tan \phi /a$ at the $T$ rows.
In Cartesian coordinates they are 1, 1 and 0.
The area of a cell is $A^t = \cos \phi^t \Delta x^t \Delta y^t$ with $\Delta x$ and $\Delta y$ in metres.

The attributes \verb|dx| and \verb|dy| are set to None after the grid is built.
Code from wim\_simple that still uses them therefore stops with an error.

The grid has no treatment of the poles and no reduced grid.
The domain must therefore end before the poles.
The zonal width $\cos \phi\, \Delta x^t$ of a cell becomes small at high latitude.
The Courant number in $x$ then grows with $1/\cos \phi$, and the time step has to be chosen accordingly.

\subsection*{Propagation}

On the sphere the action equation of wim\_simple becomes
\beq
 \pt A + \frac{1}{a \cos \phi} \lk \p_\lambda (\dot x A) + \p_\phi (\cos \phi\, \dot y A) \rk
 + \p_\theta (\dot \theta A) + \p_\nu (\dot \nu A) = S
\eeq
with $\lambda$ and $\phi$ in radians.
The velocities $\dot x$ and $\dot y$ are the eastward and northward velocities in m/s.
They are those of wim\_simple, $\dot{\v x} = (c/2) \v s + \v U + R \v s$.
The direction $\theta$ is counted from north, with $\v s = (\sin \theta, \cos \theta)$ as in wim\_simple.

The discrete form is the flux form of wim\_simple.
The flux in $y$ is multiplied by $\cos \phi^u$.
The flux divergence is divided by $\cos \phi^t \Delta x^t$ and $\cos \phi^t \Delta y^t$, with the widths in metres.
The fluxes in frequency and direction are those of wim\_simple.
The Courant numbers of the Superbee scheme are $|\dot x| \Delta t/(\cos \phi^t \Delta x^t)$ and $|\cos \phi^u \dot y| \Delta t /(\cos \phi^t \Delta y^t)$.

The refraction by the Doppler velocity $\v U$ needs the horizontal gradients of $\v U$.
On the sphere these are $(a \cos \phi)^{-1} \p_\lambda$ and $a^{-1} \p_\phi$.
The discrete differences of wim\_simple are therefore divided by $\cos \phi^t \Delta x^u$ and $\Delta y^u$ instead of $\Delta x$ and $\Delta y$.

On the sphere a ray that keeps its direction relative to north does not follow a great circle.
The direction therefore changes along the ray even without currents.
With currents the frequency changes as well.
Both effects add terms to $\dot \theta$ and $\dot \sigma$ of wim\_simple,
\beq
 \dot \theta &=& \dot \theta_{simple} + \lk c/2 + \v s \cdot \v U \rk \sin \theta \frac{\tan \phi}{a} \\
 \dot \sigma &=& \dot \sigma_{simple} - \frac{\sigma}{2} \lk \v U \cdot \v s_\perp \rk \sin \theta \frac{\tan \phi}{a}
\eeq
with $\v s_\perp = (\cos \theta, -\sin \theta)$ and $\dot \nu = \dot \sigma /(2 \pi)$.
Here $\dot \theta_{simple}$ and $\dot \sigma_{simple}$ are the expressions of wim\_simple with the gradients on the sphere.
The added term of $\dot \theta$ is evaluated at the faces in direction, as $\dot \theta$ of wim\_simple.
The added term of $\dot \sigma$ is evaluated at the centres in direction.

Without currents only the term $(c/2) \sin \theta \tan \phi /a$ remains.
It turns the ray onto the great circle.
The frequency is then constant along the ray.

The two terms follow exactly from the dispersion relation in spherical coordinates.
The derivation is checked symbolically in \verb|development/derive_spherical_rays.py|.
The documentation of ECMWF (IFS Part VII, Eq.\ 2.25c) keeps only the term $(c/2) \sin \theta \tan \phi/a$.
The terms with $\v U$ are of the order $|\v U| \tan \phi /a$.
This is about $2 \cdot 10^{-7}\,$s$^{-1}$ at $60^\circ$ latitude for $|\v U| = 1\,$m/s.

The option \verb|limit_fdot| of wim\_simple also bounds the added terms.

The energy fixer is that of wim\_simple.
The defect $D$ of wim\_simple is here the integral over the area of the cells.
The terms of the faces in $x$ are multiplied by $\Delta y^t$.
The terms of the faces in $y$ are multiplied by $\Delta x^t$, since the flux already carries $\cos \phi^u$.
The terms of the faces in frequency and direction are multiplied by $A^t$.
On a uniform Cartesian grid this reduces to the fixer of wim\_simple.

The open coasts are those of wim\_simple.
The outflow through a face in $x$ is weighted with $\Delta y^t$ and through a face in $y$ with $\Delta x^t$.

The Fortran kernels of the propagation (\verb|propagation_backend| $=$ 'fortran') are those of wim\_simple.
The metric factors and the added terms above are handed to them.
They give the same result as the array form.

\subsection*{Coupling to an ocean with varying cell thickness}

The coupling of wim\_simple assumes ocean cells of constant thickness $\Delta z$.
pe\_realistic has cells of varying thickness $\Delta z^t$.
wim\_realistic therefore accepts $\Delta z^t$ as an array in the coupling.

The Stokes velocity, the stress and the Doppler velocity of wim\_simple use the kernel $Q = 2 k e^{2 k z}$ integrated over each ocean cell,
\beeq
 \int_{z_b}^{z_t} Q\, dz = e^{2 k z_t} - e^{2 k z_b}
\eeeq
with the top face $z_t$ and the bottom face $z_b$ of the cell.
For an array $\Delta z^t$ the faces are the partial sums of $\Delta z^t$ from the surface downward.
The array is given from the bottom to the surface, as in pe\_realistic.
The top face of a cell is then exactly the bottom face of the cell above.
The sum of the kernel over the column is therefore exactly $1 - e^{-2 k H}$ for a column of depth $H$.
With \verb|closed_bottom| the lowest cell also takes the rest $e^{-2 k H}$, so that the sum is exactly 1, as in wim\_simple.

\verb|calc_stokes| and \verb|calc_stress_work| divide the integrated kernel by the thickness $\Delta z^t$ of each cell.
The Doppler velocities need no division and are those of wim\_simple.
With a number $\Delta z$ the arithmetic is that of wim\_simple, bit for bit.

\verb|calc_phioc_work| of wim\_simple is not generalised.
It works only with a number $\Delta z$.

\subsection*{Wind input of Janssen and dissipation of WAM4}

The attribute \verb|input_scheme| selects the source terms.
The default 'wamdi' uses the input and dissipation of wim\_simple.
The value 'janssen' uses the wind input of Janssen (1991) and the dissipation of WAM4 (Bidlot et al.\ 2007).
These are the source terms of ecWAM for \verb|IPHYS| $=0$.
They are ported from the routines \verb|sinput_jan|, \verb|sdissip_jan|, \verb|stresso|, \verb|tau_phi_hf|, \verb|z0wave|, \verb|fkmean|, \verb|femeanws| and \verb|frcutindex| of ecWAM.
The nonlinear transfer is that of wim\_simple in both cases.

The friction velocity $u_*$ and its direction $\theta_u$ are prescribed, as in wim\_simple.
ecWAM computes $u_*$ from the wind speed and the roughness.
This part of ecWAM is not ported.

The input is $S_{in} = \gamma_{in} E$ with the growth rate
\beq
 \gamma_{in} = \frac{\beta_{max}}{\kappa^2} \frac{\rho_a}{\rho_w}\, \sigma\, x^2 \mu \lk \ln \mu \rk^4
\eeq
Here $x = (u_*/c + z_\alpha) \cos(\theta - \theta_u)$ and $\ln \mu = \ln(k z_0) + \kappa / x$.
$\kappa = 0.4$ is the constant of von K\'arm\'an.
The input is zero where $\cos(\theta - \theta_u) \le 0.01$ or $\ln \mu \ge 0$.
The cells with input define the wind sea used below.

The roughness length depends on the momentum that the waves take from the wind,
\beq
 z_0 = \frac{\alpha}{g} \frac{u_*^3}{\sqrt{u_*^2 - \tau_w}}
\eeq
with the constant $\alpha$ of Charnock.
The momentum flux $\tau_w$ into the waves is divided by the density of air $\rho_a$.
It is the input of momentum over all frequency bins plus the analytic tail described below,
\beeq
 \tau_w = \frac{\rho_w}{\rho_a} \left| \int \gamma_{in} E\, \sigma\, \v s\, d\theta\, d\nu + \v \tau_{tail} \right|
\eeeq
It is limited to $u_*^2/(1+\epsilon_1)$.
Without waves $\tau_w = 0$ and $z_0 = \alpha u_*^2/g$ is the roughness of Charnock.

$z_0$ and $\tau_w$ depend on each other.
In each source step $z_0$, $\gamma_{in}$ and $\tau_w$ are therefore computed \verb|janssen_ncall| times in turn, as ecWAM does.
The first pass starts from $\tau_w$ of the previous step.

The dissipation is $S_{out} = \gamma_d E$ with
\beq
 \gamma_d = c_{dis}\, 2 \pi \nu_1\, E_{tot}^2\, \bar k^4\, X \lk (1-\delta) + \delta X \rk ~,~~ X = k/\bar k
\eeq
It does not depend on direction.
$E_{tot}$ is the total energy, $\nu_1$ the mean frequency weighted with $\nu$ and $\bar k$ the mean wavenumber from the moment with $\sqrt k$.

The means are sums over the frequency bins plus an analytic tail $E_N (\nu_N/\nu)^5$ beyond the last bin $\nu_N$,
\beeq
 E_{tot} &=& \sum E_d \Delta \nu + \nu_N E_N /4 \\
 \nu_1 &=& \lk \sum \nu E_d \Delta \nu + \nu_N^2 E_N/3 \rk / E_{tot} \\
 \bar k &=& \lk \lk \sum \sqrt{k}\, E_d \Delta \nu + \frac{2 \pi}{\sqrt g} \nu_N^2 E_N/3 \rk / E_{tot} \rk^2
\eeeq
with $E_d = \int E\, d\theta$ and $E_N$ of the last bin.

The prognostic range ends at $2.5 \max(\bar \nu, \bar \nu_{ws})$, at the nearest bin in $\ln \nu$.
Here $\bar \nu$ is the mean frequency weighted with $1/\nu$ and $\bar \nu_{ws}$ the same mean over the wind sea.
The lowest frequency bin is always prognostic.
The term of Pierson and Moskowitz in the cut-off of wim\_simple is not used, as in ecWAM.

Beyond the prognostic range the spectrum is $E_c (\nu_c/\nu)^5$, as in wim\_simple.
The stress and the energy flux of this tail are computed as in \verb|tau_phi_hf| of ecWAM.
The integral over the tail uses the rule of Simpson with \verb|jtot_tauhf| $=19$ points in $\ln(\omega \sqrt{z_0/g})$.
It runs from $\max(\omega_c, x_0 g/u_*)$ to $\omega = \sqrt{g/z_0}$.
$x_0$ is the value of $u_*/c$ below which the growth rate vanishes for the roughness of Charnock.
The growth in the tail is the one along the wind times $\cos^2(\theta - \theta_u)$ for $\cos(\theta - \theta_u) > 0$.

The tail of $\tau_w$ starts at the last frequency bin of the model.
The frequency ghost bin is still zero when the input is computed.
The tails of the stress and of $\Phi_{oc}$ with \verb|tail_flux| $=$ 'input' start at the ghost bin, as in wim\_simple.
\verb|calc_stress_work| puts the tail stress into the ocean cells point by point of the rule of Simpson, with the kernel $2 k_j e^{2 k_j z}$ of each point $j$.

The port differs from ecWAM in three points.
The widths of the first and last frequency bins are those of wim\_simple, while ecWAM uses half bins there.
This changes the means by $0.05$ to $0.5\,\%$.
The tail of $\tau_w$ starts at the last model bin, as stated above.
The tail stress is booked along each direction $\v s(\theta)$.
Its component along the wind is the stress of ecWAM, and ecWAM drops the component across the wind.

The following parts of ecWAM are not ported.
These are the gustiness, the negative input, the normalisation of the growth rate, the roughness of gravity-capillary waves, the sheltering and the viscous dissipation.
The viscous term is in the code with $c_{disvis} = 0$, as in ecWAM for \verb|IPHYS| $=0$.

$\tau_w$ and $z_0$ are carried from step to step.
They are not part of the restart of the model.
A restart starts with $\tau_w = 0$, i.e.\ with the roughness of Charnock.
The iteration of $z_0$ and $\tau_w$ corrects this within the first step.

The port is compared with the routines of ecWAM compiled with f2py.
The test is \verb|janssen_ref_test.py| in \verb|development/|.
All quantities agree to round-off.

\section*{Model parameters and options}

The parameters and options of wim\_simple are also those of wim\_realistic.
The tables list only the differences.

\begin{longtable}{@{}llp{7.6cm}@{}}
\hline
Attribute & Default & Meaning \\
\hline
\endhead
\verb|dxt|, \verb|dyt| & required & widths of the cells, a number or an array without halo.
  Degrees with \verb|coord_degree|, metres otherwise.
  Replace \verb|dx| and \verb|dy| of wim\_simple, which are still accepted \\
\verb|coord_degree| & off & longitude and latitude in degrees instead of Cartesian coordinates \\
\verb|x_origin|, \verb|y_origin| & $0$ & position of the first interior $U$ point \\
\verb|radius| & $6370\,$km & radius of the earth (m) \\
\verb|input_scheme| & 'wamdi' & source terms. 'wamdi' as in wim\_simple, 'janssen' the input of Janssen and the dissipation of WAM4 \\
\hline
\end{longtable}

\noindent
With \verb|input_scheme| $=$ 'janssen' the following constants are read.
Their defaults are those of ecWAM for \verb|IPHYS| $=0$.

\begin{longtable}{@{}llp{8.6cm}@{}}
\hline
Attribute & Default & Meaning \\
\hline
\endhead
\verb|betamax| & $1.20$ & $\beta_{max}$ of the input \\
\verb|zalp| & $0.008$ & $z_\alpha$ of the input \\
\verb|alpha_charnock| & $0.0065$ & constant $\alpha$ of Charnock \\
\verb|kappa| & $0.40$ & constant of von K\'arm\'an \\
\verb|cdis| & $-1.33$ & $c_{dis}$ of the dissipation \\
\verb|delta_sdis| & $0.5$ & $\delta$ of the dissipation \\
\verb|cdisvis| & $0$ & viscous dissipation \\
\verb|rnu| & $1.5 \cdot 10^{-6}$ & kinematic viscosity (m$^2$/s) of the viscous dissipation \\
\verb|eps1| & $10^{-5}$ & $\epsilon_1$ in the limit of $\tau_w$ \\
\verb|tailfactor| & $2.5$ & factor of the prognostic cut-off \\
\verb|janssen_ncall| & $2$ & number of passes of $z_0$, $\gamma_{in}$ and $\tau_w$ per step \\
\verb|jtot_tauhf| & $19$ & number of points of the tail integral \\
\verb|rho_a| & $1.225$ & density of air (kg/m$^3$) \\
\verb|rho_w| & $1000$ & density of water (kg/m$^3$) \\
\hline
\end{longtable}

\noindent
The densities \verb|rho_a| and \verb|rho_w| act only in the source terms of Janssen.
The input of wim\_simple and the tails of the coupling use $1.225/1000$.

\subsection*{Snapshot diagnostics}

The snapshot classes are those of wim\_simple.
They write the coordinates $x^t$ and $y^t$ of the grid, in degrees with \verb|coord_degree|.
The fields are those of wim\_simple.

\subsection*{Tests}

The tests are in \verb|pyOM3/wim_realistic/development/|.
\verb|reduction_test.py| compares wim\_realistic on a uniform Cartesian grid with wim\_simple.
The action agrees bit for bit, and to $10^{-15}$ with the energy fixer.
\verb|sphere_test.py| tests the propagation on the sphere.
The total action is conserved to $10^{-16}$ over three days.
A wave packet follows the great circle to $0.15^\circ$ in position and $0.6^\circ$ in direction after $36^\circ$ of arc.
The error of a uniform state decreases with the second power of the grid spacing.
The absolute energy closes to $2 \cdot 10^{-16}$ with a sheared current.
\verb|janssen_ref_test.py| compares the source terms of Janssen with ecWAM, see above.

\subsection*{Limits}

\begin{itemize}
\item The domain must end before the poles.
\item The widths of the cells must vary smoothly, also across a periodic seam.
\item The terms with $\v U$ in the turning on the sphere are not tested by a dynamic test.
\item \verb|calc_phioc_work| works only with a number $\Delta z$.
\item With \verb|input_scheme| $=$ 'janssen' and \verb|tail_flux| $=$ 'wamdi' the tails of the stress use the dissipation of wim\_simple.
\item $\tau_w$ and $z_0$ of the scheme of Janssen are not part of the restart.
\end{itemize}

\end{document}
